\section*{Asesmen}
\addcontentsline{toc}{section}{Asesmen}

Bab X memuat Praktikum 2 (8\%) dan bagian UTS (20\%). Fokusnya adalah
FIRST/FOLLOW, LL(1), dan implementasi parser.

\begin{table}[H]
\centering
\caption{Instrumen asesmen Bab X}
\label{tab:ch10-asesmen}
\begin{tabularx}{\textwidth}{@{}lXl@{}}
\toprule
\textbf{Instrumen} & \textbf{Indikator} & \textbf{Bobot} \\
\midrule
Praktikum 2 & Parser, AST, 20 tes, diagnostik & 8\% \\
UTS & FIRST/FOLLOW, tabel LL(1), penelusuran & 20\% \\
Kuis & Recursive descent dan pemulihan & formatif \\
\bottomrule
\end{tabularx}
\end{table}

Kriteria kelulusan: parser menerima program valid, menolak program
tidak valid dengan pesan yang informatif, dan membangun AST yang
sesuai dengan struktur program.
